\section{Fisher information rank deficiencies of the pre-G1 model: derivations}\label{app:fim}

This appendix supplies the derivations underlying the identifiability
analysis of Section~\ref{sec:fim}. We derive analytically the three
rank-deficient parameter pairs of the pre-G1 FSA-v2 model
(Sections~\ref{app:fim-pair1}--\ref{app:fim-pair3}), describe the
resulting null space of the FIM (Section~\ref{app:fim-null}), and
prove that the G1 reparametrization is mathematically equivalent to
the pre-G1 model under a fixed relabelling of parameters
(Section~\ref{app:fim-equiv}).

Throughout, we use the operating point of
Definition~\ref{def:operating-point}, write
$\delta A = A - A_{\mathrm{typ}}$ and $\delta F = F - F_{\mathrm{typ}}$
for state fluctuations, and recall that at one-day windows the slow
Banister state moves by at most $O(\mathrm{d}t/\tau)$ where
$\tau \in \{\tau_B, \tau_F\}$, so $\delta A$ and $\delta F$ are
small in this regime.

\subsection{Pair 1: \texorpdfstring{$(\kappa_B, \varepsilon_A)$}{(kappa\_B, epsilon\_A)}}\label{app:fim-pair1}

The pre-G1 $B$-drift is
\begin{equation*}
  \dd B = \bigl[\kappa_B(1+\varepsilon_A A)\Phi - B/\tau_B\bigr]\dd t + \cdots
\end{equation*}
Linearising the factor $1+\varepsilon_A A$ around $A_{\mathrm{typ}}$
and writing $\delta A = A - A_{\mathrm{typ}}$,
\begin{equation}
  \kappa_B(1+\varepsilon_A A)\,\Phi
   \;=\;
   \underbrace{\kappa_B(1+\varepsilon_A A_{\mathrm{typ}})}_{=:\,\kappa_B^{\mathrm{eff}}}\,\Phi
   \;+\; \kappa_B\,\varepsilon_A\,\delta A\,\Phi.
  \label{eq:pair1-expand}
\end{equation}
The first term carries the dependence of the daily drift on the
parameters via the single product $\kappa_B^{\mathrm{eff}}$; the
second term is $O(\delta A)$, vanishing at the operating point.

Now consider how the data depend on small perturbations
$(\delta\kappa_B,\delta\varepsilon_A)$. Differentiating
\eqref{eq:pair1-expand} at the truth point gives
\begin{equation*}
  \frac{\partial}{\partial \kappa_B}\bigl[\kappa_B(1+\varepsilon_A A)\Phi\bigr]
   = (1+\varepsilon_A A)\,\Phi
   = (1+\varepsilon_A A_{\mathrm{typ}})\,\Phi + \varepsilon_A \delta A\,\Phi,
\end{equation*}
\begin{equation*}
  \frac{\partial}{\partial \varepsilon_A}\bigl[\kappa_B(1+\varepsilon_A A)\Phi\bigr]
   = \kappa_B A\,\Phi
   = \kappa_B A_{\mathrm{typ}}\,\Phi + \kappa_B \delta A\,\Phi.
\end{equation*}
At the operating point (where $\delta A = 0$) these two gradients
are linearly dependent: their ratio is
$(1+\varepsilon_A A_{\mathrm{typ}})/(\kappa_B A_{\mathrm{typ}})$,
independent of $\Phi$. The score function with respect to
$\kappa_B$ and the score function with respect to $\varepsilon_A$
are therefore proportional, with proportionality constant fixed by
the truth values. Their outer product in \eqref{eq:fim-def} is rank
one, not rank two.

The combination that the data \emph{can} resolve is precisely the
product $\kappa_B^{\mathrm{eff}} = \kappa_B(1+\varepsilon_A A_{\mathrm{typ}})$.
The orthogonal direction
$(\delta\kappa_B, \delta\varepsilon_A) \perp \nabla \kappa_B^{\mathrm{eff}}$
moves the parameter vector but leaves the daily drift unchanged at
the operating point.

\subsection{Pair 2: \texorpdfstring{$(\mu_F, \mu_{FF})$}{(mu\_F, mu\_FF)}}\label{app:fim-pair2}

The pre-G1 Stuart--Landau growth rate is
$\mu(B,F) = \mu_0 + \mu_B B - \mu_F F - \mu_{FF} F^2$. Completing
the square in the quadratic in $F$,
\begin{align}
  -\mu_F F - \mu_{FF} F^2
  &= -\mu_F F - \mu_{FF}\bigl[(F-F_{\mathrm{typ}})^2 + 2 F_{\mathrm{typ}} F - F_{\mathrm{typ}}^2\bigr]
       \nonumber\\
  &= -\bigl(\,\underbrace{\mu_F + 2 F_{\mathrm{typ}} \mu_{FF}}_{=:\,\mu_F^{\mathrm{eff}}}\,\bigr) F
     \;-\; \mu_{FF}(F-F_{\mathrm{typ}})^2 + \mu_{FF} F_{\mathrm{typ}}^2.
  \label{eq:pair2-expand}
\end{align}
The constant $\mu_{FF} F_{\mathrm{typ}}^2$ shifts the intercept
$\mu_0$ but is uninformative about $\mu_F$ or $\mu_{FF}$ separately;
the linear-in-$F$ slope $\mu_F^{\mathrm{eff}}$ does carry their
combined effect; and the quadratic
$-\mu_{FF}(F-F_{\mathrm{typ}})^2$ is $O(\delta F^2)$, so it
contributes at second order in the small fluctuations of $F$ at
one-day windows.

Differentiating $\mu(B,F)$ with respect to $\mu_F$ and
$\mu_{FF}$ at the truth point and setting $\delta F = 0$,
\begin{equation*}
  \partial_{\mu_F}\,\mu(B,F)\big|_{F_{\mathrm{typ}}} = -F_{\mathrm{typ}},
  \qquad
  \partial_{\mu_{FF}}\,\mu(B,F)\big|_{F_{\mathrm{typ}}} = -F_{\mathrm{typ}}^2.
\end{equation*}
Both gradients point in the same one-dimensional subspace (they
differ only by the scalar factor $F_{\mathrm{typ}}$), so the FIM
contribution from $(\mu_F,\mu_{FF})$ at the operating point is rank
one. The combination resolved is
$\mu_F^{\mathrm{eff}} = \mu_F + 2 F_{\mathrm{typ}} \mu_{FF}$ (the
slope at $F_{\mathrm{typ}}$); the orthogonal direction is silent.

\subsection{Pair 3: \texorpdfstring{$(\tau_F, \lambda_A)$}{(tau\_F, lambda\_A)}}\label{app:fim-pair3}

The pre-G1 $F$-clearance rate is
$(1+\lambda_A A)/\tau_F$. Linearising in $\delta A$,
\begin{equation}
  \frac{1+\lambda_A A}{\tau_F}
   \;=\;
   \frac{1+\lambda_A A_{\mathrm{typ}}}{\tau_F}
   + \frac{\lambda_A}{\tau_F}\,\delta A
   \;=:\;
   \frac{1}{\tau_F^{\mathrm{eff}}} + O(\delta A),
  \qquad
  \tau_F^{\mathrm{eff}} := \frac{\tau_F}{1+\lambda_A A_{\mathrm{typ}}}.
  \label{eq:pair3-expand}
\end{equation}
The leading term is the inverse effective acute timescale, set by
the product $(1+\lambda_A A_{\mathrm{typ}})/\tau_F$; the $O(\delta A)$
correction vanishes at the operating point. By the same argument as
in Pair~1, the gradients of the $F$-decay rate with respect to
$\tau_F$ and $\lambda_A$ at the operating point are scalar multiples
of one another, so the corresponding FIM contribution is rank one.
The combination resolved is the inverse effective rate
$1/\tau_F^{\mathrm{eff}}$.

\subsection{The null space of the FIM}\label{app:fim-null}

Combining the three rank-one defects of
Sections~\ref{app:fim-pair1}--\ref{app:fim-pair3}, the FIM of the
pre-G1 model evaluated at the operating point has a null space of
dimension three. The three null vectors point along
\begin{equation}
  (\delta\kappa_B,\delta\varepsilon_A) \perp \nabla \kappa_B^{\mathrm{eff}},
  \quad
  (\delta\mu_F,\delta\mu_{FF}) \perp \nabla \mu_F^{\mathrm{eff}},
  \quad
  (\delta\tau_F,\delta\lambda_A) \perp \nabla(1/\tau_F^{\mathrm{eff}}),
  \label{eq:fim-nulls}
\end{equation}
i.e.\ along perturbations of each pair that hold the strongly
identified linear combination fixed. The total FIM rank is
therefore $d - 3$, where $d$ is the full dimension of the parameter
vector.

The remaining $d - 3$ identified directions include all observation-model
parameters (which couple linearly to the latent state and are
therefore generically identifiable from a one-day record under the
sleep- and wake-gating structure of
\eqref{eq:obs-hr}--\eqref{eq:obs-steps}), the slow Banister
parameters $\tau_B$ and $\kappa_F$ (which appear linearly and not in
collinear pairs), and the diffusion scales $\sigma_B,\sigma_F,\sigma_A$
(which are identifiable from the empirical quadratic variation of
the observation residuals).

\subsection{The G1 reparametrization is mathematically equivalent to the pre-G1 model}\label{app:fim-equiv}

\begin{proposition}[G1 is mathematically equivalent to v2]\label{prop:g1-equiv}
The G1 drift \eqref{eq:fsa-B}--\eqref{eq:fsa-A} with the truth values
of Table~\ref{tab:fsa-truth} is identically equal, as a function of
$(B,F,A,\Phi)$, to the pre-G1 drift under the relabelling
\begin{align*}
  \kappa_B^{\mathrm{old}} &= \kappa_B^{\mathrm{new}}/(1+\varepsilon_A A_{\mathrm{typ}}), &
  \tau_F^{\mathrm{old}}   &= \tau_F^{\mathrm{new}}\,(1+\lambda_A A_{\mathrm{typ}}), \\
  \mu_F^{\mathrm{old}}    &= \mu_F^{\mathrm{new}} - 2 F_{\mathrm{typ}} \mu_{FF}, &
  \mu_0^{\mathrm{old}}    &= \mu_0^{\mathrm{new}} - \mu_{FF} F_{\mathrm{typ}}^2,
\end{align*}
with $\varepsilon_A,\lambda_A,\mu_{FF},\mu_B,\eta,\tau_B,\kappa_F$
all unchanged.
\end{proposition}

\begin{proof}
Substitute the relabelling into the pre-G1 drift terms
\eqref{eq:fsa-B}--\eqref{eq:fsa-A} (in their original v2 form) and
collect:

\emph{$B$-equation.} The pre-G1 $B$-drift is
$\kappa_B^{\mathrm{old}}(1+\varepsilon_A A)\Phi - B/\tau_B$.
Substituting
$\kappa_B^{\mathrm{old}} = \kappa_B^{\mathrm{new}}/(1+\varepsilon_A A_{\mathrm{typ}})$
gives
$\kappa_B^{\mathrm{new}} \cdot (1+\varepsilon_A A)/(1+\varepsilon_A A_{\mathrm{typ}}) \cdot \Phi - B/\tau_B$,
which is exactly the G1 $B$-drift \eqref{eq:fsa-B}.

\emph{$F$-equation.} The pre-G1 $F$-drift is
$\kappa_F\Phi - (1+\lambda_A A) F/\tau_F^{\mathrm{old}}$. Substituting
$\tau_F^{\mathrm{old}} = \tau_F^{\mathrm{new}}(1+\lambda_A A_{\mathrm{typ}})$
gives
$\kappa_F\Phi - (1+\lambda_A A)/(1+\lambda_A A_{\mathrm{typ}}) \cdot F/\tau_F^{\mathrm{new}}$,
which is exactly the G1 $F$-drift \eqref{eq:fsa-F}.

\emph{$A$-equation.} The growth rate $\mu(B,F)$ becomes, under the
relabelling,
\begin{align*}
  \mu_0^{\mathrm{old}} + \mu_B B - \mu_F^{\mathrm{old}} F - \mu_{FF} F^2
  &= \bigl(\mu_0^{\mathrm{new}} - \mu_{FF} F_{\mathrm{typ}}^2\bigr) + \mu_B B
     - \bigl(\mu_F^{\mathrm{new}} - 2 F_{\mathrm{typ}} \mu_{FF}\bigr) F - \mu_{FF} F^2 \\
  &= \mu_0^{\mathrm{new}} + \mu_B B - \mu_F^{\mathrm{new}} F
     - \mu_{FF}\bigl(F^2 - 2 F_{\mathrm{typ}} F + F_{\mathrm{typ}}^2\bigr) \\
  &= \mu_0^{\mathrm{new}} + \mu_B B - \mu_F^{\mathrm{new}} F
     - \mu_{FF}(F - F_{\mathrm{typ}})^2,
\end{align*}
which is exactly the G1 growth rate \eqref{eq:fsa-mu}. The $A$-drift
is then $\mu(B,F) A - \eta A^3$, identical in both
parametrizations. \qed
\end{proof}

The numerical truth values of Table~\ref{tab:fsa-truth} satisfy the
above relabelling exactly, so the unit test
\texttt{tests/test\_g1\_reparam.py} -- which evaluates the two
drift formulas at $50$ random states -- agrees to floating-point
precision. The reparametrization is purely a coordinate change in
parameter space; the dynamics, the observation model, and the joint
likelihood $p(y_{1:T} \mid \theta)$ are unchanged.
